\section{Experimental design}
\label{sec:experiments}
\label{sec:exp-setup}

The experiments ask whether retaining a shared nuisance variable helps a learned predictor combine incomplete sensor readings, and whether that benefit survives changes in sensor count, calibration size and data source. These are different questions. A model can transfer to a larger sensor set while still fitting real response functions poorly. We therefore retain the synthetic comparisons, the losses on real networks, the component ablations and the execution audits as separate pieces of evidence. All numerical results in this manuscript come from the existing research record. No model was fitted or selected when preparing this version.

\subsection{Synthetic tasks and evaluation panels}

For each task, the response of sensor $j$ is generated as
\begin{equation}
 x_j=a_jy+b_j+\alpha\sin(f_jy)+d_ju+\sigma_j\epsilon_j,
 \label{eq:experimental-generator}
\end{equation}
where the target $y$, shared nuisance $u$ and sensor specific noise $\epsilon_j$ are independent standard normal variables. The loading $a_j$ has a random sign and magnitude drawn uniformly from $[0.6,1.8]$, $b_j\sim\mathcal{N}(0,0.4^2)$, $f_j\sim\mathcal{U}[0.5,2.5]$, $d_j\sim\mathcal{N}(0,0.7^2)$ and $\sigma_j\sim\mathcal{U}[0.15,0.8]$. Parameters are fixed within an episode and redrawn between tasks. The default nonlinearity is $\alpha=0.4$. Sharing $u$ makes the sensor errors dependent even when the target is known, while the sinusoidal term introduces a controlled departure from the linear Gaussian assumptions of the anchor.

Each default episode contains 48 labeled support rows and 48 query rows. Every support sensor entry is removed independently with probability 0.2. Input centering and scaling use observed support entries, and query labels are used only for scoring. During synthetic training, a query has either no missing sensor or one missing sensor. Evaluation considers $k\in\{0,1,2,3\}$ missing sensors. For a given $k$, every deletion pattern is scored when there are at most 20 patterns; otherwise, 20 patterns are sampled with a fixed seed. Query rows and masks are averaged within each task before computing the task level comparison. Consequently, hundreds of scores from the same task do not become hundreds of independent observations.

The main protocol v2 panel, F1, contains 256 new tasks with five sensors. Six additional panels contain 128 tasks each and change one factor at a time: eight sensors, sixteen sensors, $\alpha=0.8$, $\alpha=0$, 16 support rows or 96 support rows. The model trained on five sensors is used without adaptation on the wider panels. Protocol v3 adds G1 with 256 new tasks and five sensors, and G3 with 128 new tasks and sixteen sensors. These panels test the transformer readout hypothesis formed after the v2 results; reusing F1 for this hypothesis would have given only descriptive evidence.

\subsection{Comparators and privileged references}

The closed form comparators are fitted from the episode support only. FA Gaussian retains one residual factor, while EM Gaussian fits a joint Gaussian to the incomplete support using expectation maximization~\citep{dempster1977em}. The reported closed form fits use 40 EM iterations. The neural anchor uses eight iterations, so initialization at the anchor does not make the untrained neural model identical to the stronger reported FA Gaussian baseline. On F1 with two missing sensors, the latter scores \FfaOnekc\ and the untrained network scores \Untrainedkc. This distinction matters when attributing an improvement to learning.

NL FA fits cubic response curves in the standardized target, with an unpenalized intercept and slope, curvature penalty $\lambda=100$, and a factor model for the residual covariance. Inference uses a 401 point grid. Curvature is clamped outside the support target range to obtain linear extrapolation. The polynomial degree and penalty were selected on development data; unrestricted selection by leaving out support rows produced unstable extrapolation. Bayesian linear regression uses evidence maximization and includes parameter uncertainty in its predictive variance~\citep{mackay1992bayesian}. Complete case ridge and the earlier ridge using mean imputation are retained to show the effect of strengthening the original control. When too few complete rows are available, the implementations use their recorded fallback rules, so the complete case labels describe the preferred fitting rule rather than an unconditional guarantee. The Gaussian process and nearest neighbour comparisons were added after the protocols and are descriptive.

The learned controls include a residual MLP with 15,490 parameters, the earlier scalar cavity architecture, a transformer with 202,562 parameters, and a transformer trained on every query mask. The residual MLP receives the same FA anchor and support encodings as the small lifted model, together with readings, masks and a support summary. Its sensor inputs are concatenated, so it is defined only for five sensors. The transformer follows the cell token design of TabPFN v2~\citep{hollmann2025tabpfnv2}, but is trained for this study from initialization. It is not the pretrained TabPFN v2 model. Four blocks alternate attention across columns and across support rows, using width 64 and four heads. Random feature embeddings allow it to operate at different sensor counts. Protocol v3 uses the same backbone with a lifted information site readout, giving 202,822 parameters.

We also evaluate two privileged references. The oracle knows the true generating parameters and evaluates the posterior numerically on a dense target grid. A second reference uses Hamiltonian Monte Carlo with the generator priors~\citep{neal2011hmc}. It integrates over response parameters and analytically marginalizes the shared nuisance, but fixes each response sign from its support correlation. Finite noisy support can identify that sign incorrectly, so this is a sign conditioned approximation rather than a certified Bayes optimum. It uses four chains, 50 retained draws per chain and query specific reweighting. Reported convergence diagnostics have $\widehat{R}<1.1$ for at least 99.9\% of identifiable parameters on the evaluated panels. Agreement with the oracle when true parameters are supplied checks the predictive calculation; it does not establish exact posterior exploration. Any percentage of a gap closed is relative to this particular HMC reference.

\subsection{Training and real data episodes}

The small learned models use a pool of 60,000 synthetic tasks generated with seed 5150. Training runs for 20,000 AdamW steps, with batches of 128 tasks and eight query rows per task, learning rate $1.8\times10^{-3}$, cosine decay to one tenth of that rate, weight decay $10^{-4}$ and gradient clipping at 5. The principal lifted model and stated controls use three training seeds. The transformers instead see newly generated tasks at every update and run for 30,000 steps with 16 tasks and 16 queries per task. They use learning rate $10^{-3}$, 1,000 warmup steps, cosine decay and clipping at 1. The v3 comparison fixes the same recipe, seed and final checkpoint for both transformer variants. These comparisons match relevant inputs and masks, but do not match total task exposure or every training resource across all architectures. They therefore evaluate the recorded recipes rather than a universal parameter efficiency law.

Beijing multi site air quality~\citep{zhang2017beijing} provides twelve monitoring stations. Each episode predicts a station's log pollutant concentration from contemporaneous readings at the other eleven stations. A week is eligible when the target station has at least 96 usable hours; 48 support hours and 48 distinct query hours are sampled within that week. Sensor gaps remain as observed in the raw data. Source adaptation uses March 2013 through December 2015, and evaluation uses January 2016 through February 2017. Protocol v2 evaluates PM$_{2.5}$ on \Rbeijingepisodes\ station and week episodes. Protocol v3 evaluates previously unscored NO$_2$ and CO targets on \VthreebeijingnoTwoepisodes\ and \Vthreebeijingcoepisodes\ episodes. Each target uses its own eligible episodes. Conditions add either three or six sensor deletions to the natural query masks, while retaining an observed sensor whenever one exists. Support and query hours are sampled from the same week. This measures contemporaneous interpolation after local calibration, not prediction of a future time period from past hours alone.

The UCI Air Quality dataset~\citep{devito2008airquality} uses five metal oxide channels on one device to predict log CO or log NO$_2$ from reference analyzers. The loader removes the documented invalid readings and rows without the required sensor and target values before constructing episodes. Missing individual channels are then simulated, including 20\% support dropout. The fitting period ends before October 1, 2004; the retained later test windows start on October 6 and comprise 23 weeks. CO was examined during development, and NO$_2$ was the v1 confirmation target. Their later reuse is secondary evidence on the same device, not replication across independent datasets.

Real data fine tuning uses one fixed recipe: 4,000 source episodes, 2,000 updates, learning rate $3\times10^{-4}$ and random seed 11. Synthetic initialization, source only initialization and specified component controls are compared with the same adaptation rule. The gas sensor drift dataset~\citep{vergara2012drift} was inspected during development but excluded before building a test split: extreme readings made Gaussian mean NLL depend on a few episodes, with mean and median differing by more than 100 nats. This excluded benchmark remains a limitation of domain coverage.

\subsection{Scores, intervals and declared decisions}

Negative log predictive density in nats is the primary score. For Gaussian outputs with mean $\mu$ and variance $v$, it is $\frac12\{\log(2\pi v)+(y-\mu)^2/v\}$. Lower values are better; negative NLL is possible for continuous densities. Secondary scores are coverage of the central 90\% interval, continuous ranked probability score and RMSE~\citep{gneiting2007scoring}. Likelihood scores depend on the target scale and transformation, so values for different pollutants are not an absolute common accuracy scale.

Gains subtract the candidate NLL from its comparator NLL. A positive gain favors the candidate. Normal approximation intervals use the mean paired gain plus or minus 1.96 standard errors. Synthetic uncertainty is computed over tasks after averaging masks and the stated fitted seeds. Beijing differences are averaged across target stations within each week before computing the interval across weeks. This avoids counting stations sharing a week as independent, but does not establish independence of adjacent weeks. Averaging the NLL of fitted seeds is also distinct from scoring an ensemble density, and the intervals condition on those fits rather than incorporating all training randomness. The declared superiority gates generally require a gain of at least 0.01 nats and a positive lower bound. E7b instead requires a lower bound greater than $-0.02$ nats. These are nominal intervals for individual endpoints, with no multiplicity adjustment or declared overall superiority rule.
